% ============================
%   1. Introduction
% ============================
\section{Introduction}
\label{sec:introduction}

\subsection{Context and motivation}

Network attacks have moved from a niche concern to a mainstream business
risk: modern enterprises rely on continuously inspecting their traffic for
malicious patterns, and the discipline that formalises that inspection is
\textit{Network Intrusion Detection}. Two complementary schools of thought
dominate the field. \textbf{Signature-based} systems --- Snort,
Suricata and Zeek being the canonical references~\cite{roesch1999snort,
suricata, paxson1999bro} --- pattern-match every packet (or every flow)
against a hand-curated rule base, producing high-precision alerts but
inevitably missing attacks for which no rule has yet been
written~\cite{garcia2009anomaly}. \textbf{Anomaly-based} (or
machine-learning-based) systems instead learn a model of ``normal''
behaviour and flag deviations~\cite{buczak2016survey}; they cope better
with previously-unseen attacks but pay for that flexibility with a higher
false-positive rate and the well-known interpretability deficit famously
critiqued by Sommer and Paxson~\cite{sommer2010outside}.

Real-world Security Operations Centres rarely deploy one school in
isolation. They run a signature engine for the cheap, high-confidence
detections \textit{and} a machine-learning detector as a second opinion,
and they need a way to investigate any alert in detail---often by
replaying the precise packets that led up to it. That last requirement
(\textit{forensics}) is the most under-served in open-source academic
work: the typical research prototype evaluates a classifier on a static
dataset and stops there, leaving the analyst experience entirely out of
scope.

\subsection{Problem statement}

The project addresses three concrete questions:

\begin{enumerate}
    \item Can a single self-hostable application combine signature and
        machine-learning detection in a way where the two engines reinforce
        each other (\textit{hybrid} detection) rather than running
        side-by-side as silos?
    \item Can such a system persist enough context per alert that an
        analyst can perform meaningful triage without an external packet
        sniffer? Specifically: which packets immediately preceded the
        alert, what was their payload, what was the inter-arrival
        timing?
    \item Can the whole thing be packaged so that a non-developer (a
        coordinator, a defense committee, an evaluator at the
        \textit{Faculty of Mathematics and Computer Science}) can launch
        it with a double-click and see results within seconds, without
        Python, Node or virtual-environment knowledge?
\end{enumerate}

\subsection{Contribution of this stage}

At the time of this progress report, the project contributes:
\begin{itemize}
    \item A modular detection pipeline (\code{Dispatcher} -- see
        Section~\ref{sec:architecture}) where signature and ML verdicts
        are merged into three explicit alert categories:
        \textit{signature-only}, \textit{ml-only}, and \textit{hybrid}.
    \item A YAML-seeded, hot-reloadable rule store editable from the UI,
        avoiding the restart-to-edit penalty common in Snort/Suricata
        deployments.
    \item A per-alert ring-buffered packet capture written as JSON Lines
        on disk, downloadable from the UI through an authenticated
        endpoint.
    \item A React + Tailwind SIEM console (login, dashboard, alerts,
        rules, forensics, ML models, settings) with live WebSocket
        updates for both packet rate and alert stream.
    \item A one-click cross-platform launcher
        (\code{run.py}~/~\code{run.bat}~/~\code{run.sh}) that creates a
        virtual environment, installs dependencies, starts both
        services, opens the browser, and handles graceful
        shutdown---designed specifically for the
        \textit{open-and-demonstrate} scenario typical of a thesis
        defense.
\end{itemize}

The remainder of this document is organised as follows.
Section~\ref{sec:background} reviews the necessary background.
Section~\ref{sec:related_work} situates the project relative to existing
tools and research. Section~\ref{sec:architecture} describes the system
design. Section~\ref{sec:implementation} discusses the implementation
details. Section~\ref{sec:current_state} reports what is functional
today, with quantitative observations gathered during development.
Section~\ref{sec:future_work} closes with the items remaining before the
final defense.
